\section{Introduction}
\label{sec:intro}

Eye tracking enables direct measurement of human visual attention and is
widely used in robotics, augmented reality, transportation, healthcare,
and gaming. Two hardware paradigms dominate. Screen-based trackers offer
sub-degree accuracy at the cost of a fixed setup; head-mounted trackers
(e.g.\ Pupil Labs Core) are more flexible but rely heavily on per-session
calibration to map raw pupil estimates to screen coordinates. The
calibration step thus determines the practical accuracy of any
downstream task, from clinical assessment to interactive games.

Classical calibration methods fit a parametric form (similarity transform,
polynomial regression, radial basis function (RBF), or hybrid
similarity-plus-RBF) to a small set of calibration fixations collected at
the start of each session. These methods routinely achieve $0.6$--$1.3^\circ$
on commodity head-mounted hardware, but their accuracy degrades when the
calibration set is small, when head pose drifts during the session, or
when corneal geometry varies across users. A recent line of work
\emph{neural refinement} therefore proposes to combine the predictions of
several classical baselines through a lightweight neural network, with
the hope that the network learns to weigh them adaptively.

\paragraph{What we found.}
We set out to extend one such pipeline (originally targeting ICML 2026)
and to address reviewer concerns about subject-disjoint evaluation. In
the course of reproducing its headline numbers we discovered that the
publicly available implementation contains a single line that subtracts
the regression target from each baseline's input feature
(\verb|mb_features = baseline_pred - target|). The network's input
therefore contains the very quantity it is asked to predict. The
per-split z-score normalization that the code applies on top conceals
the leakage at test time but does not eliminate it; the network exploits
the residual structure to drive training error very close to zero, and
test error settles near a value that is dominated by the per-split
normalization statistics rather than by genuine generalization. Removing
this single line:
\begin{itemize}
\item changes the JuDo1000 headline from $5.8$~px to $30.5$~px
  (which is \emph{worse} than the $21.7$~px similarity baseline),
\item changes the self-collected dataset headline from $42.3$~px (claimed
  $0.96^\circ$) to $\sim 45$~px (statistically indistinguishable from the
  $45.5$~px sim+RBF baseline), and
\item leaves the proposed design with no measurable advantage over a
  carefully fit classical calibrator.
\end{itemize}
Figure~\ref{fig:leakage} summarizes the diagnosis.

\paragraph{What works.}
The same dataset analysis exposes a different opportunity. Per-trial
biases of $25$--$60$~px are present and direction-stable
(Figure~\ref{fig:dataset}c), suggesting that a brief
\emph{online recalibration} after the standard calibration phase could
identify and remove most of the residual error. Indeed, simply
subtracting the empirical mean residual computed from $K=12$ extra
fixations reduces mean L2 by $15\%$. However, the empirical mean is
noisy at small $K$. We propose a tiny MLP ($\sim$1\,000 parameters) that
takes per-trial bias statistics (mean magnitude, variance, $K$) and
predicts a 2-D shrinkage factor $\lambda \in [0,1]^2$ to apply to the
mean. Trained leave-one-subject-out across our twelve participants and
evaluated on the held-out subject, the learned shrinkage estimator
beats both the raw mean and a globally-tuned fixed $\lambda$ at every
$K \geq 5$, achieving a $17\%$ reduction in mean L2 at $K=12$
(Figure~\ref{fig:shrinkage}).

\paragraph{Contributions.}
\begin{itemize}
\item \textbf{An audit of multi-baseline neural refinement} for
  head-mounted eye tracking. We pinpoint a single-line input-feature
  leak that drives the published gains, and we re-evaluate six neural
  designs (anchored MLP residual, softmax-blend over baselines,
  fourier-position-encoded MLP, DeepSets context encoder,
  per-trial Bayesian model averaging, leakage-free top-$M$ selection)
  under an honest protocol on two datasets. None reliably beats the
  strongest classical baseline.
\item \textbf{Honest per-trial shrinkage estimation.} A
  $\sim$1\,000-parameter MLP that maps per-trial bias statistics to a
  2-D shrinkage coefficient, trained leave-one-subject-out, delivers
  a $17\%$ reduction in mean L2 over the sim+RBF anchor at $K=12$
  context fixations.
\item \textbf{Open release.} We release leakage-free reproducible
  pipelines, target-disjoint JuDo1000 splits, leave-one-subject-out
  protocol, and per-experiment results to facilitate honest
  benchmarking of future neural calibration methods.
\end{itemize}

\begin{figure}[t]
\centering
\includegraphics[width=\linewidth]{fig02_leakage.pdf}
\caption{Honest evaluation eliminates the published headline numbers.
We replace \texttt{mb\_features = baseline\_pred - target} with raw
baseline predictions normalized by training-set statistics. (a) On
JuDo1000, the honest pipeline yields $30.5$~px --- worse than the
$21.7$~px similarity baseline. The leaky pipeline reproduces the
published $5.8$~px. (b) On our 12-subject dataset the same bug inflates
training-time signal; the honest neural method matches the strongest
classical baseline rather than improving over it.}
\label{fig:leakage}
\end{figure}
